\origpage{361}{N. Robertson, P.D. Seymour / Journal of Combinatorial Theory, Series B 90 (2004) 325--385}
\noindent is a circuit $F'$ with $E(F') \subseteq E(F) \cup \{e\}$, $ins(F) \subseteq ins(F')$, and $|E(F')| < |E(F)|$. But then $(\mathscr{F} \setminus \{F\}) \cup \{F'\}$ is a feast of smaller order than $\mathscr{F}$, a contradiction.

The result follows. $\square$

\textbf{11.4.} \textit{Let $\mathscr{F}$ be a feast of minimum order $\le (m+1)n$ and subject to that with $\Sigma(\lambda(F) : F \in \mathscr{F})$ minimum. Then each $F \in \mathscr{F}$ is $\lambda$-linked to $\mathscr{T}$.}

\textbf{Proof.} Let $F \in \mathscr{F}$, and choose $F'$ as in 8.5. Then $r^* \notin ins(F') \setminus U(F')$, because $|E(F')| \le |E(F)| \le (m+1)n$ and $(\Gamma, \phi, \lambda, \mathscr{T})$ is not $(m+1)n$-flawed in freedom. Thus $r^* \in V(F')$, and since $ins(F) \subseteq ins(F')$ it follows that $\mathscr{F}' = (\mathscr{F} \setminus \{F\}) \cup \{F'\}$ is a feast. From the minimality of the order of $\mathscr{F}$ we deduce that $|E(F)| = |E(F')|$; but then $\lambda(F') \le \lambda(F)$, and again it follows that $\lambda(F') = \lambda(F)$. Hence $F' = F$, and so $F$ is $\lambda$-linked to $\mathscr{T}$, as required. $\square$

If $(\Gamma, \phi, \lambda, \mathscr{T}) \in \mathscr{D}$ is $(n, (\omega_1, \ldots, \omega_m))$-flawed on $S$, it follows from 11.2, 11.3, 11.4 that there is a disjoint feast $\mathscr{F}$ of order $\le (m+1)n$, such that each of its members is $\lambda$-linked to $\mathscr{T}$. We call $\mathscr{F}$ an $((\omega_1, \ldots, \omega_m), S)$-flaw for $(\Gamma, \phi, \lambda, \mathscr{T})$.

\textbf{Proof of 11.1.} As usual, it suffices to prove that if $\mathscr{D}' \subseteq \mathscr{D}$, and each member of $\mathscr{D}'$ has the same $((\omega_1, \ldots, \omega_m), S)$-flaw $\mathscr{F}$, then there is a well-behaved set $\mathscr{C}$ of $\chi$-places, such that for all $(\Gamma, \phi, \lambda, \mathscr{T}) \in \mathscr{D}'$ there is a rooted location which $(m+1)n$-isolates $\mathscr{T}$ and for which $(\Gamma, \phi, \mathscr{L}) \in \mathscr{C}$. Let the components of $S \setminus I(\mathscr{F})$ be $X_1, \ldots, X_t$. For each $(\Gamma, \phi, \lambda, \mathscr{T}) \in \mathscr{D}'$ there exist $\omega'_1, \ldots, \omega'_t \in \{\omega_1, \ldots, \omega_p\}$, such that for $1 \le i \le t$, $\phi(e) \not\ge \omega'_i$ for each edge $e$ of $\Gamma$ with $e \subseteq X_i$. Since there are only finitely many such $t$-tuples $(\omega'_1, \ldots, \omega'_t)$, we may further assume that the same $t$-tuple $(\omega'_1, \ldots, \omega'_t)$ works for all members of $\mathscr{D}'$. For $1 \le i \le t$, let $\Omega_i$ be the ideal of $\Omega_\chi(S)$ with element set $\{x \in E(\Omega_\chi(S)) : x \not\ge \omega'_i\}$. Then for $1 \le i \le t$, $\Omega_i < \Omega_\chi(S)$, and for all $(\Gamma, \phi, \lambda, \mathscr{T}) \in \mathscr{D}'$, $\phi(e) \in E(\Omega_i)$ for all $e \in E(\Gamma)$ with $e \subseteq X_i$.

Let $(\Gamma, \phi, \lambda, \mathscr{T}) \in \mathscr{D}'$. For each $F \in \mathscr{F}$, let $\Delta(F) = ins(F)$; this does not depend on the choice of $(\Gamma, \phi, \lambda, \mathscr{T})$, since $|\mathscr{F}| \ge 2$ and $ins(F)$ meets no member of $\mathscr{F}$ except $F$. Since $(\Gamma, \phi, \lambda, \mathscr{T})$ is not $((m+1)n+1)$-flawed in distance and $dist(\Gamma) > \frac{1}{2}((m+1)n+1)$, it follows that $F$ meets no cuff except $\bar{r}$.

Let $\Sigma'$ be obtained from $\Sigma_\chi$ by deleting $\Sigma_\chi \cap (\Delta(F) \setminus U(F))$ for each $F \in \mathscr{F}$. Then $\Sigma'$ is homeomorphic to one of the fragments obtained by cutting $\Sigma_\chi$ along $U(\bigcup(F : F \in \mathscr{F}))$, and to simplify notation we assume that $\Sigma'$ is such a fragment, and the associated surjection $\psi$ is the identity when restricted to $\Sigma'$. We see that $X_1, \ldots, X_t$ are long sides of $\psi^{-1}(\Phi_\chi) \cap \Sigma'$. Let $\chi'$ be the colour scheme with $\Sigma_{\chi'} = \Sigma'$, $\Phi_{\chi'} = \psi^{-1}(\Phi_\chi) \cap \Sigma'$, $\Omega_{\chi'}(k) = \Omega_\chi(k)$ $(k = 2, 3)$, $\Omega_{\chi'}(X_i) = \Omega_i$ $(1 \le i \le t)$ and $\Omega_{\chi'}(S') = \Omega_\chi(R)$ for each long side $S' \ne X_1, \ldots, X_t$ of $\Phi_{\chi'}$, where $R \ne S$ is the long side of $\Phi_\chi$ with $S' \subseteq R$. Let us define $f(X_i) = S$ $(1 \le i \le t)$, and $f(S') = R$ for each long side $S'$ of $\Phi_{\chi'}$ with $S' \ne X_1, \ldots, X_t$, where $R$ is the long side of $\Phi_\chi$ with $S' \subseteq R$. Then $f$ is an embedding of $\chi'$ in $\chi$, and so since $\chi$ satisfies $S_4$, $\chi'$ is not orientedly bad. But now the result follows as in the proofs of 10.2, 10.4, 10.5. $\square$
